% =====================================================================
% chapters/01-introduction.tex
% PT: Capítulo 1 -- Introdução. Exemplo de: parágrafo de abertura,
%     secções e subsecções, acrónimos (\gls), citação (\cite), nota de
%     rodapé, listas itemize/enumerate e referências cruzadas a capítulos
%     e secções (\cref, \ref).
% EN: Chapter 1 -- Introduction. Shows: opening paragraph, sections and
%     subsections, acronyms (\gls), citation (\cite), footnote,
%     itemize/enumerate lists and cross-references to chapters and
%     sections (\cref, \ref).
% =====================================================================

% PT: Cada capítulo tem um \label para poder ser referido noutros sítios.
%     Prefixos sugeridos: chap: capítulo, sec: secção, fig: figura,
%     tab: tabela, eq: equação, app: apêndice.
% EN: Each chapter has a \label so it can be referred to elsewhere.
%     Suggested prefixes: chap: chapter, sec: section, fig: figure,
%     tab: table, eq: equation, app: appendix.
\chapter{Introduction}
\label{chap:introduction}

% PT: Parágrafo de abertura: em poucas frases, diz ao leitor o que o
%     capítulo contém, antes da primeira secção.
% EN: Opening paragraph: in a few sentences, tell the reader what the
%     chapter contains, before the first section.
This chapter shows how a chapter is written with this template. In a real
dissertation, the introduction presents the context of the work, the problem
it addresses, its objectives and research questions, and the structure of the
document. Every chapter of the template contains a short English example
followed by a Portuguese example; together they cover the features you are
most likely to need. Replace the text and keep the structure.

% ---------------------------------------------------------------------
% PT: Exemplo em inglês
% EN: English example
% ---------------------------------------------------------------------
\section{English example}
\label{sec:intro:en}

\subsection{Context and Motivation}
\label{sec:intro:context}

% PT: \gls escreve a sigla por extenso no primeiro uso e abreviada depois.
% EN: \gls spells out the acronym on first use and abbreviates it afterwards.
Research in \gls{ai}, and in particular in \gls{ml}, has changed the way
data are analysed in many fields \cite{LeCun2015}. From the second use
onwards, only the short form is printed: \gls{ai} and \gls{ml}.
% PT: Nota de rodapé: o texto vai dentro de \footnote{...}.
% EN: Footnote: the text goes inside \footnote{...}.
The context section should explain why the topic matters and who benefits
from the work.\footnote{Footnotes are numbered per chapter. Use them sparingly,
for side remarks that would interrupt the main text.}

\subsection{Objectives}
\label{sec:intro:objectives}

The general objective of a dissertation is usually stated in one sentence and
then broken down into specific objectives:

% PT: Lista com marcas (itemize).
% EN: Bulleted list (itemize).
\begin{itemize}
  \item review the state of the art on the topic;
  \item design and implement the proposed solution;
  \item evaluate the solution against a baseline;
  \item discuss the results and their limitations.
\end{itemize}

\subsection{Research Questions}
\label{sec:intro:questions}

% PT: Lista numerada (enumerate). O enumitem permite rótulos próprios.
% EN: Numbered list (enumerate). enumitem allows custom labels.
\begin{enumerate}[label=\textbf{RQ\arabic*.}, leftmargin=*]
  \item Which methods are currently used to solve the problem?
  \item How does the proposed solution compare with those methods?
  \item Under which conditions does the proposed solution fail?
\end{enumerate}

\subsection{Document Structure}
\label{sec:intro:structure}

% PT: \cref escreve o tipo e o número ("Chapter 2"). \Cref é a forma
%     para o início da frase; com a opção capitalise, usada neste
%     template, as duas dão o mesmo resultado em inglês.
% EN: \cref writes the type and the number ("Chapter 2"). \Cref is the
%     form for the start of a sentence; with the capitalise option, used in
%     this template, both give the same result in English.
\Cref{chap:review} reviews the related work. \Cref{chap:methodology}
describes the methodology and \cref{chap:results} reports the results, which
are discussed in \cref{chap:discussion}. \Cref{chap:conclusion} concludes the
document, and \cref{app:example} holds supporting material. The objectives
listed in \cref{sec:intro:objectives} are revisited in the conclusion.

% ---------------------------------------------------------------------
% PT: Exemplo em português
% EN: Portuguese example
% ---------------------------------------------------------------------
\section{Exemplo em português}
\label{sec:intro:pt}

% PT: O ambiente otherlanguage aplica a hifenização portuguesa ao bloco.
% EN: The otherlanguage environment applies Portuguese hyphenation to the block.
\begin{otherlanguage}{portuguese}

\subsection{Enquadramento}

% PT: Sigla portuguesa definida em frontmatter/acronyms.tex.
% EN: Portuguese acronym defined in frontmatter/acronyms.tex.
A \gls{ia} tem vindo a ser aplicada em áreas tão diferentes como a
indústria, a saúde e a educação \cite{Goodfellow2016}. Depois do primeiro
uso, a sigla aparece só abreviada: \gls{ia}. O enquadramento deve explicar
a relevância do tema e o problema concreto que a dissertação
resolve.\footnote{As notas de rodapé funcionam da mesma forma em texto
português.}

\subsection{Objetivos}

% PT: Em texto português, escreva o nome à mão antes de \ref.
% EN: In Portuguese text, write the name by hand before \ref.
Os objetivos específicos seguem a ordem dos capítulos, descrita na
Secção~\ref{sec:intro:structure}:

\begin{enumerate}
  \item rever o estado da arte (Capítulo~\ref{chap:review});
  \item conceber e implementar a solução (Capítulo~\ref{chap:methodology});
  \item avaliar e discutir os resultados (Capítulos~\ref{chap:results}
        e~\ref{chap:discussion}).
\end{enumerate}

\end{otherlanguage}
